% ============================================================
%  COURSE ID – Course Title
%  Assignment / Report Template
% ============================================================
\documentclass[12pt, a4paper]{article}

% ── Encoding & fonts ────────────────────────────────────────
\usepackage[T1]{fontenc}
\usepackage[utf8]{inputenc}
\usepackage{lmodern}

% ── Page layout ─────────────────────────────────────────────
\usepackage[
  top=2.5cm, bottom=2.5cm,
  left=2.5cm, right=2.5cm
]{geometry}

% ── Mathematics ─────────────────────────────────────────────
\usepackage{amsmath, amssymb, amsthm}

% ── Graphics & tables ───────────────────────────────────────
\usepackage{graphicx}
\usepackage{float}
\usepackage{booktabs}
\usepackage{multirow}
\usepackage{array}

% ── Code listings ───────────────────────────────────────────
\usepackage{listings}
\usepackage{xcolor}

\lstset{
  backgroundcolor=\color{gray!10},
  basicstyle=\ttfamily\small,
  keywordstyle=\color{blue}\bfseries,
  commentstyle=\color{green!50!black}\itshape,
  stringstyle=\color{red!70!black},
  numberstyle=\tiny\color{gray},
  numbers=left,
  breaklines=true,
  frame=single,
  captionpos=b,
  showstringspaces=false
}

% ── Hyperlinks ──────────────────────────────────────────────
\usepackage[colorlinks=true, linkcolor=blue, citecolor=blue, urlcolor=blue]{hyperref}

% ── Headers & footers ───────────────────────────────────────
\usepackage{fancyhdr}
\pagestyle{fancy}
\fancyhf{}
\rhead{COURSE ID – Report 1}
\lhead{University Name}
\cfoot{\thepage}

% ── Paragraph spacing ───────────────────────────────────────
\usepackage{parskip}
\setlength{\parskip}{6pt}

% ── Caption formatting ──────────────────────────────────────
\usepackage[font=small, labelfont=bf]{caption}

% ── Bibliography ────────────────────────────────────────────
\usepackage[numbers, sort&compress]{natbib}

% ============================================================
%  Title information  ← EDIT HERE
% ============================================================
\title{
  \vspace{-1cm}
  {\large University Name} \\[0.3cm]
  {\LARGE \textbf{COURSE ID: Course Title}} \\[0.4cm]
  {\Large Report 1} \\[3cm]
  
}

\author{
  \textbf{Student Name:} Your Full Name \\
  \textbf{Student ID:}   XXXXXXXXX \\
  \textbf{Email:}        yourname@university.edu
}

\date{Submission Date: \today}

% ============================================================
\begin{document}

\maketitle
\thispagestyle{empty}
\newpage

% ── Abstract ────────────────────────────────────────────────
\begin{abstract}
Write a concise summary (150–250 words) of your report here.
Briefly describe the objectives, the methods used, the main
results, and your conclusions.
\end{abstract}

\newpage
\tableofcontents
\newpage

% ============================================================
\section{Introduction}
\label{sec:intro}

Introduce the topic and background of the assignment. State the
objectives clearly. You may reference related work
\cite{example2020} here.

% ============================================================
\section{Background \& Related Work}
\label{sec:background}

Discuss any relevant theory, prior work, or datasets used.

% ============================================================
\section{Methodology}
\label{sec:method}

Describe your approach, model architecture, or experimental
setup.

\subsection{Dataset}

Describe the dataset used (size, preprocessing steps, etc.).

\subsection{Model / Approach}

Explain your model or algorithmic approach. You can include
equations:

\begin{equation}
  \mathcal{L} = -\sum_{i=1}^{N} y_i \log(\hat{y}_i)
  \label{eq:crossentropy}
\end{equation}

\subsection{Experimental Setup}

Detail hyperparameters, hardware, training procedure, etc.

% ============================================================
\section{Results}
\label{sec:results}

Present your experimental results. Use tables and figures to
support your analysis.

\subsection{Quantitative Results}

\begin{table}[H]
  \centering
  
  \label{tab:results}
  \begin{tabular}{lccc}
    \toprule
    \textbf{Model} & \textbf{Accuracy} & \textbf{F1} & \textbf{Loss} \\
    \midrule
    Baseline        & 0.75              & 0.74        & 0.62          \\
    Our Model       & \textbf{0.87}     & \textbf{0.86} & \textbf{0.41} \\
    \bottomrule
  \end{tabular}
  \caption{Comparison of model performance.}
  \label{tab:results}
\end{table}

\subsection{Figures}

% Uncomment and replace with your actual figure file:
% \begin{figure}[H]
%   \centering
%   \includegraphics[width=0.8\linewidth]{figures/training_curve.png}
%   \caption{Training and validation loss over epochs.}
%   \label{fig:training}
% \end{figure}

% ============================================================
\section{Discussion}
\label{sec:discussion}

Analyse and interpret your results. Compare against baselines.
Discuss limitations and potential improvements.

% ============================================================
\section{Conclusion}
\label{sec:conclusion}

Summarise your findings and their significance. Suggest future
directions for research.

% ============================================================
%  Code snippet example (optional)
% ============================================================
\appendix
\section{Code Snippets}

\begin{lstlisting}[language=Python, caption={Example training loop.}]
for epoch in range(num_epochs):
    model.train()
    for batch in train_loader:
        optimizer.zero_grad()
        outputs = model(**batch)
        loss = outputs.loss
        loss.backward()
        optimizer.step()
    print(f"Epoch {epoch+1}, Loss: {loss.item():.4f}")
\end{lstlisting}

% ============================================================
%  Bibliography
% ============================================================
\bibliographystyle{plainnat}
\bibliography{references}

% If you prefer inline references, use thebibliography:
% \begin{thebibliography}{9}
% \bibitem{example2020}
%   Author, A.\ (2020).
%   \textit{Title of the paper}.
%   Journal Name, 10(2), 100--115.
% \end{thebibliography}

\end{document}
